\chapter{Upgrading}\label{ch:upgrade}

\emph{You want to install new versions of Agern safely, and go back if one misbehaves.} The
\cmd{agern} command downloads a version from GitHub, verifies its checksum, installs it with
\cmd{apt-get}, restarts Agern and waits until it answers with the new version. If it does not come
up, the previous version is reinstalled automatically.

\section{Which version you run}\label{sec:version}

\begin{enumerate}
  \item In the web UI, open the user menu (your name, top right). Its last line shows the running
    version (figure~\ref{fig:user-menu}).
  \item On the server, \cmd{agern version} shows the installed version, the previous one kept for a
    rollback, and what the running server reports.
\end{enumerate}

\screenshot[0.45\linewidth]{48-user-menu-desktop}{The user menu with the running version (here the demo's).}{user-menu}

There are two \term[channel]{channels}:

\begin{itemize}
  \item \textbf{stable} (the default): tagged releases such as \texttt{1.4.0}. Releases with a suffix,
    such as \texttt{1.5.0-rc.1}, are pre-releases and are only installed when you name them.
  \item \textbf{main}: the rolling build \texttt{main-latest}, rebuilt after every change merged into
    Agern's \texttt{main} branch. Its versions look like \texttt{1.4.1-main.20261119101500.g1a2b3c4}:
    after the last release, before the next, in commit order. Choose it only for a factory whose
    users accept that a bad merge reaches them minutes later.
\end{itemize}

\section{Upgrade by hand}\label{sec:upgrade}

\begin{enumerate}
  \item Upgrade to the latest release:
\begin{lstlisting}
sudo agern upgrade
\end{lstlisting}
    Or to a specific release, or to the newest build of \texttt{main}:
\begin{lstlisting}
sudo agern upgrade --version 1.4.0
sudo agern upgrade --channel main
\end{lstlisting}
  \item If runs or agents are working, Agern refuses, because a restart fails runs in progress:
    \msg{agern: error: not restarting now, the server is busy: 1 run(s) in progress, 2 agent(s)
    working. Try again later, or use -{}-force (interrupted runs fail)}. Wait, or add
    \cmd{--force} if you accept that.
  \item Check the result with \cmd{agern version}, or in the user menu.
\end{enumerate}

Your configuration in \file{/etc/agern/agern.env} is never touched by an upgrade. Database changes
run when the new version starts.

\section{Roll back}\label{sec:rollback}

\begin{enumerate}
  \item If the new version misbehaves, reinstall the previous package, which \cmd{agern upgrade}
    kept in \file{/var/cache/agern}:
\begin{lstlisting}
sudo agern rollback
\end{lstlisting}
\end{enumerate}

A rollback does not undo database changes. Rolling back across a change of the database schema is
only safe if the older version tolerates the newer schema.

\section{Upgrade automatically}\label{sec:auto-upgrade}

Agern can install new versions by itself, when nobody's work is in progress.

\begin{enumerate}
  \item Turn it on, optionally for the main channel:
\begin{lstlisting}
sudo agern auto-upgrade on
sudo agern auto-upgrade on --channel main
\end{lstlisting}
  \item Every 15 minutes a timer looks for a newer version. It installs it as soon as no run is
    queued, specifying, building, reviewing or approving and no agent is starting or mid-turn. Runs
    waiting for a human and idle chats do not count.
  \item See what it does, and turn it off again:
\begin{lstlisting}
agern auto-upgrade status
sudo agern auto-upgrade off
\end{lstlisting}
\end{enumerate}

If the factory stays busy for 24 hours (\texttt{AGERN\_UPGRADE\_MAX\_WAIT\_HOURS}), the timer gives up
on that version with a warning in the journal; \cmd{sudo agern upgrade} installs it whenever it
suits you. Details:
\repodoc{INSTALL-LINUX.md\#8-upgrades-rollback-and-automatic-upgrades}{\texttt{INSTALL-LINUX.md}, section~8}.

\whatyouhave
You know which version runs, how to upgrade by hand or automatically, from releases or from
\texttt{main}, and how to roll back. Next, protect your data: chapter~\ref{ch:backup}.
